% =====================================================================
% 01_hypothesis.tex  --  Hypothesis (Exp. 2 & 3)
% =====================================================================

\section{Hypothesis}
\label{sec:cdmech-hypothesis}

Contrastive decoding rests on a specific mechanistic story: a degraded amateur
branch amplifies prior-driven hallucinations, and subtracting its logits from
the expert steers generation toward better-grounded tokens. This story implies
two properties that should be directly observable at decoding time. We state
them as hypotheses and design Experiments~2 and~3 to measure each one on two
model families (LLaVA-1.5-7B and Qwen2.5-VL-7B).

\begin{description}
  \item[\textbf{H1 (Intervention).}] Contrastive decoding should change the
  emitted token at the steps where hallucinations arise. Experiment~2 measures
  how often the contrastive output coincides with the expert's own greedy
  choice, and compares this rate at real- versus hallucinated-object steps, to
  quantify where and how often the intervention takes effect.
  \item[\textbf{H2 (New grounding).}] When contrastive decoding does intervene,
  it should introduce a better-grounded alternative --- changing which
  candidates are available, rather than re-ordering a shortlist the expert has
  already fixed. Experiment~3 measures the overlap between the expert, amateur,
  and contrastive candidate sets, to determine whether the contrast supplies new
  candidates or re-weights existing ones.
\end{description}

\noindent Together these experiments characterize the contrastive step
mechanistically: Experiment~2 establishes \emph{when} it acts, and Experiment~3
establishes \emph{what} it has available to act on. This complements the
external, benchmark-level analysis of \citet{yin2025mirage} by examining the
decoding process itself, on the generative captioning setting measured by CHAIR.
